\section{Introdução}
O desenvolvimento de programas no mundo atual depende frequentemente do uso de
APIs, muitas vezes desenvolvidas por equipes completamente desconexas.
Frequentemente, a utilização de funções provenientes de tais fontes espera que
uma lógica interna seja seguida para evitar erros como \textit{memory leaks}, 
arquivos que permanecem abertos sem necessidade ou até \textit{use after free}s.
Tais descuidos podem levar a nenhuma consequência grave, graves erros de
segurança ou apenas comportamentos inesperados pela aplicação; as implicações de
violar os padrões de uso de uma API são muitas vezes imprevisíveis.

Contudo, descobrir quais padrões de uso são necessários por uma API é muitas
vezes difícil. A fonte principal para tal tipo de informação seria a
documentação da API em questão, porém muitas vezes esta se mostra insuficiente
ou limitada. Um exemplo claro de tal são as páginas do manual Linux para as
funções \textit{fopen} e \textit{fclose} na linguagem C \cite{man_fopen}
\cite{man_fclose}.

Qualquer desenvolvedor de C informado saberia dizer que um \textit{fclose}
precisa ser chamado após qualquer \textit{fopen}. Porém as páginas do manual de
ambas estas funções se referem a outra apenas tangencialmente em seus textos.
Desta forma, frequentemente, o aprendizado de padrões de uso se passa não por
uma simples leitura da documentação, mas sim como algo visto em exemplos de
código ou através de discussões informais com outros desenvolvedores.

Devido a tal, os reais padrões de uso que devem ser respeitados quando ao
utilizar APIs permanecem obscuros. Assim, vários estudos foram e continuam a
serem feitos para detectar tais de forma automática.

O presente estudo, então, busca aplicar técnicas desenvolvidas na área para a
detecção de padrões de uso na \textit{glibc}. Os objetivos principais são
realizar uma análise comparativa das diferentes formas que tal mineração de
dados pode ser feita, quais as vantagens e desvantagens de cada e gerar
catálogos dos padrões de uso encontrados por cada método.  
